%-*- coding:utf-8 -*-

\begin{frame}[t,fragile]{ビニングによる誤差評価}
  \begin{itemize}
    %\setlength{\itemsep}{1em}
  \item 連続する \(n\) 個の測定値を、一つのビンにまとめる
  \item 測定値 \(M\) 個を \(m=M/n\) 個のビンに分け、ビンごとの平均を計算
    \[
    \overline A_k^{(n)}
    =\frac{1}{n}\sum_{i=n(k-1)+1}^{nk}A_i
    \qquad(k=1,\ldots,m)
    \]
    ビン平均の平均は、元のデータ全体の平均 \(\overline A\) に等しい

  \item ビン平均の分散を求める
    \[
    \left(s^{(n)}\right)^2
    =\frac{1}{m-1}\sum_{k=1}^{m}
    \left(\overline A_k^{(n)}-\overline A\right)^2
    \]

  \item 全体の平均値\(\overline{A}\)の分散の推定値を求める。ビン平均の分散を、さらにビン数 \(m\) で割る
    \[
    \boxed{
    \left(\sigma^{(n)}\right)^2
    =\frac{1}{m}\left(s^{(n)}\right)^2
    =\frac{n}{M}\left(s^{(n)}\right)^2
    }
    \]

  \item ビン幅を変えて、誤差の推定値が安定することを確認する
  \end{itemize}
\end{frame}
